\honest{The Ritual}{Eliezer Yudkowsky}{2008}

This is a story. Jeffreyssai, a ``Bayesian Master,'' receives a guest in the formal room he keeps for outsiders, because it is what they expect. \nb{In the book's order this is the first appearance of Jeffreyssai and the beisutsukai. Yudkowsky had introduced them in earlier stories, such as ``The Failures of Eld Science'' (May 2008), which the book places later.}

The guest, from another Conspiracy, has told him something. ``Are you sure?'' he asks. ``Does it not suffice that I am a domain expert, and you are not?'' she answers. She has never had formal training, but she has heard the beisutsukai explain themselves often enough. Her counsel has always been good. She jokes about statistics, then admits she is not sure: only he can decide. But he should hold his ritual for when he must ``very seriously consider abandoning a long-held premise of your existence.'' An expert has told him he is probably wrong, so he concedes. She kisses his hand, says ``Farewell, sensei,'' and ``It won't be the same.'' I do not say what the premise is.

He cancels tomorrow's classes. For the rest of the day he does nothing in particular, blocking every thought he has had before and every conclusion. \nb{This is the advice of ``Crisis of Faith,'' posted the day before: rest the previous day, and ``Clear your mind of all standard arguments.''}

In the morning he goes down to a plain white room. On the wall are plaques: ``That which can be destroyed by the truth should be. People can stand what is true, for they are already enduring it. Curiosity seeks to annihilate itself.'' \nb{The first line is P. C. Hodgell's, the second Eugene Gendlin's, the third from Yudkowsky's ``Twelve Virtues of Rationality.'' The story names none of them.}

Below them is a tally. Seven times he has entered this room; five times he kept his belief, twice he came out a different person. If the plus column stayed empty, you would have to admit you lacked the ability, unless you had been born knowing everything. \nb{So a belief kept after the ritual is no failure. ``Crisis of Faith'' had counted Robert Aumann's kept faith as a failed doubt, because it took the faith to be false.}

He sits facing the blank wall and rehearses ``seven major principles and sixty-two specific techniques,'' and his ``fourteen most embarrassing oversights.'' I name none of them. Then he asks himself the question. I do not say what it is, or how it comes out. \nb{The story closes the book \textsc{How to Actually Change Your Mind}. It shows the preparation for changing one's mind, and stops before the change.}
